Para as bolsas de produtividade em Comunicação, o CNPq considera, na chamada de 2026, até trinta trabalhos no cálculo da nota de produção bibliográfica: pelo menos dez devem ser artigos em periódicos qualificados, embora livros, coletâneas e capítulos também contem (CNPq, “Critérios para Avaliação de Bolsas de Produtividade”, seção “Programa Básico de Comunicação”).
O Manifesto de Leiden lembra que as práticas de publicação e citação variam entre as áreas e que os indicadores devem auxiliar quem avalia a pesquisa, sem dispensar a leitura do trabalho (Hicks et al., “The Leiden Manifesto”, princípios 1, 6 e 7).
Saber em que coluna entraria esta carta ainda não me diria o que vim a saber ao escrevê-la.
